% 荣枯鉴
% 荣枯鉴.tex

\documentclass[12pt,UTF8]{ctexbook}

% 设置纸张信息。
\input{../../../../Book/common/page}

% 设置字体，并解决显示难检字问题。
\xeCJKsetup{AutoFallBack=true}
\setCJKfamilyfont{hei}{SimHei}
\setCJKfamilyfont{kai}{KaiTi}

% 目录 chapter 级别加点（.）。
\usepackage{titletoc}
\titlecontents{chapter}[0pt]{\vspace{3mm}\bf\addvspace{2pt}\filright}{\contentspush{\thecontentslabel\hspace{0.8em}}}{}{\titlerule*[8pt]{.}\contentspage}

% 支持目录点击跳转
\usepackage[colorlinks,linkcolor=blue,citecolor=blue,urlcolor=blue]{hyperref}

% 设置 part 和 chapter 标题格式。
\ctexset{
	part/name= {第,卷},
	part/number={\chinese{part}},
	chapter/name={第,章},
	chapter/number={\arabic{chapter}}
}

% 图片相关设置。
\usepackage{graphicx}
\graphicspath{{Images/}}

% 设置署名格式。
\newenvironment{shuming}{\hfill\zihao{4}}

% 注脚每页重新编号，避免编号过大。
\usepackage[perpage]{footmisc}

% 设置古文原文格式。
\newenvironment{yuanwen}{\bfseries\zihao{4}}

% 设置段落间距
\setlength{\parskip}{0.8em plus 0.2em minus 0.1em}

% 列表格式支持，列表项向右偏移。
\usepackage{enumitem}

\title{\heiti\zihao{0} 荣枯鉴}
\author{冯道}
\date{\today}

\begin{document}

\maketitle
\tableofcontents

\frontmatter
\chapter{前言}

《荣枯鉴》又名《小人经》，为五代时期冯道所著。冯道历仕后唐、后晋、后汉、后周四朝，事奉从唐庄宗至周世宗多达十位皇帝，始终位居宰相、三公、三师之位，是五代时期著名的政治人物。

本书以独特的视角剖析了官场处世之道，揭示了人情世故中的种种玄机。书名"荣枯"二字，取草木之繁荣与枯萎，比喻人之得势与失势、顺境与逆境。全书共分五卷，分别为圆通、闻达、解厄、交结、节义，每卷皆有其独特的处世智慧。

\mainmatter

\part{圆通卷}

\chapter{原文}

\begin{yuanwen}

善恶有名，智者不拘也。天理有常，明者不弃也。道之靡通，易者无虞也。
惜名者伤其名，惜身者全其身。名利无咎，逐之非罪，过乃人也。
君子非贵，小人非贱，贵贱莫以名世。君子无得，小人无失，得失无由心也。名者皆虚，利者惑人，人所难拒哉。
荣或为君子，枯必为小人。君子无及，小人乃众，众不可敌矣。名可易事难易也，心可易命难易也，人不患君子，何患小人焉？



名不足以致善，名可以佛恶。

舍善而趋恶，非恶之善乎？

奸有益，人皆可为奸；

忠致祸，人难为忠。

奸有益，人皆可为奸；

忠致祸，人难为忠。

忠奸难辨也，惟心乎？

惟其利乎？惟其名乎？

名者，实之宾也。

实者，名之主也。

名实相符，世所同愿也。

名实相悖，世所难免也。

是以君子务实，小人务名。

务实者，虽无名而有实；

务名者，虽无实而有名。

故名者，虚器也；

实者，真宰也。

智者重实轻名，愚者重名轻实。

是以智者常吉，愚者常凶。

吉者，实之所致也；

凶者，名之所致也。

故智者务其实，不务其名；

愚者务其名，不务其实。

实则荣之本，名乃枯之末。

务实者荣，务名者枯。

荣枯之间，惟心自知。
\end{yuanwen}

\chapter{释义}

\section{名实之辨}

善与恶都有其名声，但智者不会被名声所束缚；天理有其恒常的规律，但明智的人不会因此放弃变通。

名声不足以让人行善，但名声可以掩饰恶行。舍弃善行而趋向恶行，难道不是恶行伪装成善吗？

做奸佞之人有时能得到好处，所以人人都可能做奸佞；忠诚有时会招致祸患，所以人很难保持忠诚。

忠诚与奸佞难以分辨，到底是看内心？看利益？还是看名声？

名声是实际的附属，实际是名声的主宰。名实相符是世人共同的愿望，但名实不符也是世间难免的事。

因此君子追求实际，小人追求名声。追求实际的人，虽然没有名声但有实绩；追求名声的人，虽然没有实绩但有名声。

所以名声是虚浮的器物，实际是真正的主宰。智者重视实际轻视名声，愚者重视名声轻视实际。因此智者常常吉利，愚者常常凶险。

吉利是实际带来的，凶险是名声带来的。所以智者追求实际，不追求名声；愚者追求名声，不追求实际。

实际是繁荣的根本，名声是枯萎的末梢。追求实际的人繁荣，追求名声的人枯萎。繁荣与枯萎之间，只有自己的内心最清楚。

\part{闻达卷}

\chapter{原文}

\begin{yuanwen}

仕不计善恶，迁无论奸小。悦上者荣，悦下者蹇。君子悦下，上不惑名。小人悦上，下不惩恶。下以直为美，上以媚为忠。直而无媚，上疑也；媚而无直，下弃也。上疑祸本，下弃毁誉，荣者皆有小人之谓，盖固本而舍末也。
富贵有常，其道乃实。福祸非命，其道乃察。实不为虚名所羁，察不以奸行为耻。无羁无耻，荣之义也。
求名者莫仕，位非名也。求官者莫名，德非荣也。君子言心，小人攻心，其道不同，其效自异哉。

直无媚，上不惑也；

媚无直，下不奉也。

上以媚为忠，故媚者荣；

下以直为美，故直者蹇。

荣者，上之所悦也；

蹇者，下之所悦也。

上悦下怨，上荣下蹇；

上怨下悦，上蹇下荣。

是以君子处上，不媚上而悦下；

处下，不悦下而媚上。

处上而悦下，上不疑也；

处下而媚上，下不怨也。

上不疑则荣，下不怨则安。

荣安之道，在于闻达也。

闻者，名之所归也；

达者，利之所至也。

有名有利，闻达之谓也。

无名无利，困顿之谓也。

故闻达者，人之所欲也；

困顿者，人之所恶也。

欲荣而恶枯，人之常情也。
\end{yuanwen}

\chapter{释义}

\section{上下之道}

做官不计较善恶，升迁不论奸佞还是小人。让上司高兴的人荣耀，让下属高兴的人困顿。

君子让下属高兴，上司不会被名声迷惑；小人让上司高兴，下属不会惩罚恶行。

下属以正直为美德，上司以谄媚为忠诚。正直不谄媚，上司不迷惑；谄媚不正直，下属不奉迎。

上司以谄媚为忠诚，所以谄媚的人荣耀；下属以正直为美德，所以正直的人困顿。

荣耀是上司喜欢的人，困顿是下属喜欢的人。上司高兴下属怨恨，上司荣耀下属困顿；上司怨恨下属高兴，上司困顿下属荣耀。

因此君子处于上位，不谄媚上司而取悦下属；处于下位，不取悦下属而谄媚上司。

处于上位而取悦下属，上司不会怀疑；处于下位而谄媚上司，下属不会怨恨。

上司不怀疑就会荣耀，下属不怨恨就会安定。荣耀安定的道理，在于闻达。

闻是名声的归宿，达是利益的所在。有名声有利益，叫做闻达；没名声没利益，叫做困顿。

所以闻达是人们想要的，困顿是人们厌恶的。想要荣耀厌恶枯萎，这是人之常情。

\part{解厄卷}

\chapter{原文}

\begin{yuanwen}

无忧则患烈也。忧国者失身，忧己者安命。祸之人拒，然亦人纳；祸之人怨，然亦人遇。君子非恶，患事无休；小人不贤，畲庆弗绝。
上不离心，非小人难为；下不结怨，非君子勿论。祸于上，无辩自罪者全。祸于下，争而罪人者免。
君子不党，其祸无援也。小人利交，其利人助也。道义失之无惩，祸无解处必困，君子莫能改之，小人或可谅矣。

祸于心摄，心若无求，祸无从至。

患生于欲，欲若无望，患无从起。

是以无忧则患烈，无欲则患微。

君子不患无福，患无福之容；

不患无位，患无位之守。

福者，人之所欲也；

位者，人之所求也。

欲福而得祸，求位而得辱，

世之常也，人之不幸也。

是以君子务本，本立而道生。

本者，心也；

道者，路也。

心正路平，心邪路仄。

正邪之间，惟心自知。

故君子慎独，小人肆欲。

慎独者，心正也；

肆欲者，心邪也。

心正则吉，心邪则凶。

吉凶之兆，见于未然。

是以君子见微知著，

小人见著不知微。

知微者吉，不知微者凶。
\end{yuanwen}

\chapter{释义}

\section{忧患之道}

没有忧患意识，祸患就会来得猛烈。为国担忧的人会丧失自身，为自己担忧的人能安于天命。

祸患从内心产生，如果内心没有贪求，祸患就无从到来。忧患产生于欲望，如果欲望没有奢望，忧患就无从兴起。

因此没有忧患意识则祸患猛烈，没有欲望则祸患微小。

君子不担心没有福气，而担心没有承载福气的器量；不担心没有地位，而担心没有守住地位的品德。

福气是人们想要的，地位是人们追求的。想要福气却得到祸患，追求地位却得到屈辱，这是世间的常态，是人的不幸。

因此君子致力于根本，根本确立了道路就会产生。根本就是内心，道路就是人生之路。内心端正道路就平坦，内心邪恶道路就狭窄。

正邪之间，只有自己的心最清楚。所以君子谨慎独处，小人放纵欲望。谨慎独处的人内心端正，放纵欲望的人内心邪恶。

内心端正就吉利，内心邪恶就凶险。吉凶的征兆，在事情发生之前就能显现。因此君子见微知著，小人见著不知微。能知微的人吉利，不能知微的人凶险。

\part{交结卷}

\chapter{原文}

\begin{yuanwen}

智不拒贤，明不远恶，善恶咸用也。顺则为友，逆则为敌，敌友常易也。
贵以识人者贵，贱以养奸者贱。贵不自贵，贱不自贱，贵贱易焉。贵不贱人，贱不贵人，贵贱久焉。
人冀人愚而自明，示人以愚，其谋乃大。人忌人明而自愚，智无潜藏，其害无止。明不接愚，愚者勿长其明。智不结怨，仇者无惧其智。
君子仁交，惟忧仁不尽善。小人阴结，惟患阴不制的。君子弗胜小人，殆于此也。

人之情多矫，世之俗多伪。

矫情以惑众，虚伪以欺世。

惑众者荣，欺世者达。

是以真情见弃，伪道大行。

君子交人以诚，小人有求于人。

诚者，心之本也；

求者，利之动也。

心本无求，利至则求。

利尽则散，求止则离。

是以君子之交淡如水，

小人之交甘如醴。

淡者，久也；

甘者，暂也。

久者，诚之所致；

暂者，利之所使。

故交之以诚，则久而益亲；

交之以利，则暂而必疏。

疏者，交之败也；

亲者，交之成也。

成败之间，惟心自知。

是以君子慎交，小人泛交。

慎交者，择人而友；

泛交者，不择而友。

择人者，知其善恶；

不择人者，不知其善恶。

知善恶者吉，不知善恶者凶。
\end{yuanwen}

\chapter{释义}

\section{交友之道}

人的情感大多矫饰，世间的习俗大多虚伪。矫饰情感来迷惑众人，虚伪做作来欺骗世人。迷惑众人的人荣耀，欺骗世人的人显达。

因此真情被抛弃，伪道大行其道。君子与人交往出于真诚，小人对人有所求。真诚是内心的根本，有所求是利益的驱动。

内心本来无所求，利益到来才有所求。利益耗尽就散去，所求停止就分离。

因此君子之交淡如水，小人之交甘如醴。淡才能长久，甜只是暂时。长久是真诚所致，暂时是利益所使。

所以以真诚相交，就会越久越亲密；以利益相交，就会暂时而必然疏远。疏远是交往的失败，亲密是交往的成功。

成败之间，只有自己的心最清楚。因此君子谨慎交友，小人广泛交友。谨慎交友的人选择人做朋友，广泛交友的人不加选择地做朋友。

选择人的人知道对方善恶，不加选择的人不知道对方善恶。知道善恶的人吉利，不知道善恶的人凶险。

\part{节义（仪）卷}

\chapter{原文}

\begin{yuanwen}

外君子而内小人者，真小人也。外小人而内君子者，真君子也。德高者不矜，义重者轻害。
人慕君子，行则小人，君子难为也。人怨小人，实则忘义，小人无羁也。难为获寡，无羁利丰，是以人皆小人也。
位高节低，人贱义薄。君子不堪辱其志，小人不堪坏其身。君子避于乱也，小人达于朝堂。
节不抵金，人困难为君子。义不抵命，势危难拒小人。不畏人言，惟计利害，此非节义之道，然生之道焉。

名者，实之宾也；

实者，名之主也。

名实相副，世之愿也；

名实相悖，世之患也。

是以君子务实，小人务名。

务实者荣，务名者枯。

荣枯之际，惟心自知。

节者，操之守也；

义者，道之从也。

守节者，不畏死；

从义者，不苟生。

不畏死者，勇也；

不苟生者，义也。

勇者，义之助；

义者，勇之主。

有勇无义，为乱；

有义无勇，为懦。

不乱不懦，惟义惟勇。

是以君子守节从义，

小人背节弃义。

守节从义者荣，

背节弃义者枯。

荣枯之分，天理昭昭。
\end{yuanwen}

\chapter{释义}

\section{节义之辨}

名声是实际的附属，实际是名声的主宰。名实相符是世人的愿望，名实不符是世间的祸患。

因此君子追求实际，小人追求名声。追求实际的人繁荣，追求名声的人枯萎。繁荣与枯萎之际，只有自己的心最清楚。

节操是品行的操守，道义是正道的依从。坚守节操的人不怕死，依从道义的人不苟且偷生。不怕死是勇，不苟生是义。

勇是义的辅助，义是勇的主宰。有勇无义就会作乱，有义无勇就会懦弱。既不作乱又不懦弱，就要既有义又有勇。

因此君子坚守节操依从道义，小人背弃节操抛弃道义。坚守节操依从道义的人繁荣，背弃节操抛弃道义的人枯萎。

繁荣与枯萎的分别，天理昭然若揭。

\part{明鉴卷}

\chapter{原文}

\begin{yuanwen}
福不察非福，祸不预必祸。福祸先知，事尽济耳。
施小信而大诈逞，窥小处而大谋定。事不可绝，言不能尽，至亲亦戒也。佯惧实忍，外恭内忌，奸人亦惑也。知戒近福，惑人远祸，俟变则存矣。
私人惟用，其利致远。天恩难测，惟财可恃。以奸治奸，奸灭自安。伏恶勿善，其患不生。
计非金者莫施，人非智者弗谋，愚者当戒哉。
\end{yuanwen}

\chapter{释义}

\section{祸福之道}

不察福分就不是真正的福气，不预先防备灾祸就一定会有灾祸。能预先知道祸福，事情就能完全成功了。

施一点小信用而大诈谋就能得逞，观察小事而大谋略就能确定。事情不能做绝，言语不能说尽，即使是至亲也要戒备。表面害怕实际残忍，外表恭敬内心猜忌，这样连奸人也会被迷惑。知道戒备就能接近福气，迷惑别人就能远离灾祸，等待时机变化就能生存。

只任用私人，利益才能长远。君主的恩宠难以预料，只有钱财可以依靠。用奸人来治理奸人，奸人消灭了自己就安全了。对邪恶的人不要讲仁慈，祸患就不会产生。

计谋没有金钱支持就不要实施，人如果不是智者就不要谋划，愚蠢的人应当引以为戒。

\part{谤言卷}

\chapter{原文}

\begin{yuanwen}
人微不诤，才庸不荐。攻其人忌，人难容也。陷其窘地人自污，谤之易也。善其仇者人莫识，谤之奇也。究其末事人未察，谤之实也。设其恶言人弗辩，谤之成也。
谤而不辩，其事自明，人恶稍减也。谤而强辩，其事反浊，人怨益增也。
失之上者，下必毁之；失之下者，上必疑之。假天责人掩私，假民言事见信，人者尽惑焉。
\end{yuanwen}

\chapter{释义}

\section{毁谤之道}

地位低微就不要争辩，才能平庸就不要推荐。攻击别人忌讳的地方，别人难以容忍。把人陷入窘境使其自取其辱，毁谤就容易成功。善待仇人让人们无法识别，这是毁谤的奇招。追究人家的细微末节人们没有察觉，这是毁谤的实效。捏造恶言让对方无法辩解，毁谤就能成功。

受到毁谤而不辩解，事情自然会明白，人们的厌恶会稍微减少。受到毁谤而强行辩解，事情反而会更浑浊，人们的怨恨会增加。

得罪了上司，下属一定会诋毁；失去了民心，上司一定会怀疑。假借天意责备别人来掩盖私心，假借民意来办事让人相信，人们都会被迷惑。

\part{示伪卷}

\chapter{原文}

\begin{yuanwen}
无伪则无真也。真不忌伪，伪不代真，忌其莫辩。
伪不足自祸，真无忌人恶。顺其上者，伪非过焉。逆其上者，真亦罪焉。求忌直也，曲之乃得。拒忌明也，婉之无失。
忠主仁也，君子仁不弃旧。仁主行也，小人行弗怀恩。君子困不惑人，小人达则背主，伪之故，非困达也。
俗礼不拘者非伪，事恶守诺者非信，物异而情易矣。
\end{yuanwen}

\chapter{释义}

\section{真伪之道}

没有虚伪就没有真实。真实不忌讳虚伪，虚伪不能代替真实，忌讳的是无法分辨。

虚伪不足以招致灾祸，真实不忌讳别人的厌恶。顺从上司的人，虚伪也不算过错。违背上司的人，真实也是罪过。请求忌讳直接，委婉才能得到。拒绝忌讳明白，婉转才能没有过失。

忠诚来源于仁爱，君子仁爱不抛弃旧人。仁爱来源于行为，小人的行为不怀念恩德。君子困顿时不迷惑别人，小人显达时就背叛主子，这是虚伪的缘故，不是困顿或显达造成的。

不拘泥于世俗礼节不算虚伪，做坏事却守信用不算守信，事物不同人情就变了。

\part{降心卷}

\chapter{原文}

\begin{yuanwen}
以智治人，智穷人背也。伏人慑心，其志无改矣。
上宠者弗明贵，上怨者休暗结。术不显则功成，谋暗用则致胜。君子制于亲，亲为质自从也。小人畏于烈，奸恒施自败也。
理不直言，谏非善辩，无嫌乃及焉。情非彰示，事不昭显，顺变乃就焉。
仁堪诛君子，义不灭小人，仁义戒滥也。恩莫弃贤者，威亦施奸恶，恩威戒偏也。
\end{yuanwen}

\chapter{释义}

\section{降心之道}

用智慧统治人，智慧穷尽人就会背叛。使人内心畏惧屈服，他的志向就不会改变了。

被上司宠爱的人不要显耀尊贵，被上司怨恨的人不要暗中勾结。权术不显露就能成功，谋略暗中使用就能取胜。君子被亲情制约，以亲情为人质自然顺从。小人害怕刚烈，奸计经常使用就会自取失败。

道理不必直说，劝谏不必善辩，没有嫌疑才能达到目的。情感不必显露，事情不必公开，顺应变化才能成功。

仁爱能诛杀君子，道义不能消灭小人，仁义要戒滥用。恩惠不要抛弃贤能的人，威严也要施加给奸恶的人，恩威要戒偏废。

\part{揣知卷}

\chapter{原文}

\begin{yuanwen}
善察者知人，善思者知心。知人不惧，知心堪御。
知不示人，示人者祸也。密而测之，人忌处解矣。君子惑于微，不惑于大。小人虑于近，不虑于远。
设疑而惑，真伪可鉴焉。附贵而缘，殃祸可避焉。结左右以观情，无不知也。置险难以绝念，无不破哉。
\end{yuanwen}

\chapter{释义}

\section{揣知之道}

善于观察的人能了解人，善于思考的人能知心。了解人就不害怕，知心就能驾驭。

知道了不要显示出来，显示出来就会有灾祸。秘密地揣测，别人的忌讳就能解除了。君子会被小事迷惑，不会被大事迷惑。小人只考虑眼前，不考虑长远。

设置疑问使对方迷惑，真伪就能鉴别了。依附权贵攀附关系，灾祸就能避免了。结交身边的人来观察情况，没有不知道的。设置险阻来断绝念头，没有不能攻破的。

\backmatter

\chapter{后记}

《荣枯鉴》一书，虽名为"小人经"，然其文辞古朴，寓意深远。冯道以其五朝为相之经历，总结官场处世之道，既有对人性深刻的洞察，也有对世事的通达之见。

读此书者，当知书中所言，乃乱世求生之术，非盛世立身之本。荣枯之间，不仅在于处世之智，更在于内心之正。愿读者以此为鉴，知荣枯之理，守正心之道。

\end{document}
